\section{Software and reproducibility}
\label{app:software}

All computations use the Python package in the repository~\cite{repo}. The numbers and tables in
this report were produced at commit {\footnotesize\href{https://github.com/uhrm/kangaroo-sequence/tree/1192116ae9ff5b3c9f32cd87c538cb3ba24d58a2}{\texttt{1192116ae9ff5b3c9f32cd87c538cb3ba24d58a2}}},
where the outputs of all scripts are stored in \texttt{analysis/output/}. The main modules are
\begin{itemize}
  \item \texttt{comma.py}: children, successors, and the block jumps of Section~\ref{sec:jumps};
  \item \texttt{child\_graph.py}: roots and complete exploration of child trees
        (Section~\ref{sec:trees});
  \item \texttt{decades.py}: the decade table, immortal entries and the earliest infinite path
        (Sections~\ref{sec:table} and~\ref{sec:consequences}).
\end{itemize}
Table~\ref{tab:reproduce} lists which script reproduces which result. Each script is run with
\texttt{uv run python} followed by the command shown. The test suite (\texttt{uv run pytest})
checks the results against~\cite{comma} and the decade table against direct computation, as
described in the validation paragraphs above.

\begin{table}[ht]
  \centering
  \footnotesize
  \begin{tabular}{@{}ll@{}}
    \toprule
    Result & Command (output in \texttt{analysis/output/}) \\
    \midrule
    Section~\ref{sec:trees}, Appendix~\ref{app:trees}
      & \texttt{-m kangaroo\_sequence.child\_graph} (\texttt{child\_trees.txt}) \\
    Section~\ref{sec:jumps} ($D_f$), Section~\ref{sec:table} ($M$, period)
      & \texttt{analysis/decade\_period.py} (\texttt{decade\_period.txt}) \\
    Section~\ref{sec:ends}, Table~\ref{tab:cr}
      & \texttt{analysis/successor\_ends.py} (\texttt{successor\_ends.txt}) \\
    Mean lifetime for $3 \le b \le 10$
      & \texttt{analysis/mean\_lifetime.py} (\texttt{mean\_lifetime.txt}) \\
    Section~\ref{sec:consequences}, Appendix~\ref{app:choices}
      & \texttt{analysis/infinite\_path.py 100} (\texttt{infinite\_path.txt}) \\
    The decade table
      & \texttt{analysis/export\_table.py} (\texttt{decade\_table\_base10.csv}) \\
    Remark~\ref{rem:index}
      & \texttt{analysis/first\_branch\_point.py} (\texttt{first\_branch\_point.txt}) \\
    Remark~\ref{rem:index}, step by step
      & \texttt{analysis/first\_branch\_point.c}, compiled with \texttt{gcc -O2} \\
    \bottomrule
  \end{tabular}
  \caption{Scripts reproducing the results of this report.}
  \label{tab:reproduce}
\end{table}

\section{The finite trees of the base-10 child graph}
\label{app:trees}

\begin{center}
  \small
  \begin{tabular}{@{}rrll@{\qquad}rrll@{}}
    \toprule
    root & leaves & height       & largest leaf     & root & leaves & height        & largest leaf    \\
    \midrule
       1 &      8 & $10^{23.29}$  & $10^{25} - 28$  &   31 &      1 & $1$           & $45$            \\
       2 &      3 & $10^{14.29}$  & $10^{16} - 82$  &   32 &      1 & $10^{7.30}$   & $10^{9} - 82$   \\
       3 &      1 & $1$           & $36$            &   37 &      1 & $195$         & $9918$          \\
       4 &      1 & $199899$      & $10^{7} - 55$   &   38 &     10 & $10^{15.32}$  & $10^{17} - 37$  \\
       5 &      4 & $10^{14.32}$  & $10^{16} - 64$  &   39 &      1 & $19$          & $918$           \\
       6 &     32 & $10^{40.31}$  & $10^{42} - 64$  &   40 &      4 & $10^{28.29}$  & $10^{30} - 82$  \\
       7 &      1 & $14$          & $936$           &   41 &      1 & $18$          & $954$           \\
       8 &      3 & $10^{11.30}$  & $10^{13} - 28$  &   42 &      1 & $1980$        & $99927$         \\
       9 &      3 & $10^{22.33}$  & $10^{24} - 19$  &   43 &      1 & $1$           & $81$            \\
      10 &      2 & $10^{11.29}$  & $10^{13} - 19$  &   49 &      1 & $192$         & $9945$          \\
      13 &     66 & $10^{85.31}$  & $10^{87} - 37$  &   50 &      3 & $10^{16.32}$  & $10^{18} - 73$  \\
      14 &      4 & $10^{9.30}$   & $10^{11} - 82$  &   51 &      8 & $10^{16.32}$  & $10^{18} - 55$  \\
      15 &      1 & $1$           & $72$            &   52 &      9 & $10^{19.29}$  & $10^{21} - 19$  \\
      16 &      4 & $10^{12.33}$  & $10^{14} - 55$  &   53 &      1 & $20$          & $972$           \\
      17 &     31 & $10^{52.32}$  & $10^{54} - 55$  &   54 &      1 & $0$           & $54$            \\
      18 &      1 & $0$           & $18$            &   62 &      2 & $205967$      & $10^{7} - 19$   \\
      19 &      1 & $21481$       & $999981$        &   63 &      1 & $0$           & $63$            \\
      20 &      - & -             & -               &   64 &    285 & $10^{130.33}$ & $10^{132} - 19$ \\
      21 &      3 & $2073$        & $99981$         &   65 &      1 & $2099$        & $99963$         \\
      25 &      1 & $192$         & $9963$          &   74 &      1 & $19$          & $945$           \\
      26 &      2 & $199508$      & $10^{7} - 28$   &   75 &      1 & $2135$        & $99954$         \\
      27 &      1 & $0$           & $27$            &   76 &      9 & $10^{22.29}$  & $10^{24} - 46$  \\
      28 &      1 & $196$         & $9927$          &   86 &      1 & $1982$        & $99918$         \\
      29 &      1 & $10^{6.31}$   & $10^{8} - 64$   &   87 &      2 & $10^{10.29}$  & $10^{12} - 82$  \\
      30 &   1008 & $10^{363.30}$ & $10^{365} - 82$ &   98 &      1 & $204$         & $9936$          \\
    \bottomrule
  \end{tabular}
\end{center}
Heights count edges from the root to the deepest leaf; height entries $10^{x}$ are rounded to two
decimals of the exponent. Every largest leaf is a landmine, given exactly.

\section{Branch choices of A367620 over one period}
\label{app:choices}

The first $314$ terms of \oeis{A399179}, the choices at the branch points of \oeis{A367620} in
decades $10, \dots, 933$ ($0$ = smaller child, $1$ = larger child); the sequence repeats them forever
(Theorem~\ref{thm:periodic}).
\begin{center}
  \small
  \texttt{0011 1001 1101 1000 0110 0101 1111 1000 0110 0100 1000 0110}\\
  \texttt{1010 0010 0001 1100 1011 0001 1001 1100 0111 1111 1001 0111}\\
  \texttt{0011 1110 0101 1111 1000 0011 0000 0000 1011 1110 1111 1010}\\
  \texttt{1110 1111 1110 1101 0011 0011 0001 0110 0100 1111 0001 1111}\\
  \texttt{0001 1011 1011 0010 0100 0110 1011 0111 0010 0101 1000 1100}\\
  \texttt{0011 1110 1110 0010 0010 1110 0001 1111 1101 0101 0000 1110}\\
  \texttt{0100 0111 1010 1110 1111 1010 11\phantom{xx xxxx xxxx xxxx xxxx xxxx}}
\end{center}
